\chapter{系统测试与效果评价}
\label{chap:evaluation}

测试是验证系统是否满足需求的必要环节。本章给出测试环境、功能与性能测试方案、问答测试集与评价指标、对照实验设计和结果分析方法。需要说明的是，本章的重点是建立可复现的测试方案和记录规范，所有指标均按系统实际运行结果计算；在系统联调完成前，不以估计值或主观判断代替实测数据。

\section{测试环境与测试方案}
\label{sec:test-env}

测试环境分为服务端、客户端和数据三部分，如表~\ref{tab:test-env} 所示。服务端使用一台开发机运行后端、MySQL 和 Chroma，嵌入模型和重排模型在本地 CPU 上加载；客户端使用主流浏览器访问前端；测试数据包括校园公开制度、办事指南、常见问答和人工编写的测试问题。为了保证结果可复现，测试前固定模型版本、提示词版本、检索参数和知识库快照，测试过程中不修改配置。

\begin{table}[H]
  \centering
  \caption{测试环境}
  \label{tab:test-env}
  \zihao{5}
  \renewcommand{\arraystretch}{1.25}
  \begin{tabular}{p{3.0cm}p{9.6cm}}
    \toprule
    类别 & 配置 \\
    \midrule
    服务端 & Windows 11，16 GB 内存，CPU 推理，Python 3.11 \\
    数据库 & MySQL 8.0，Chroma 本地持久化 \\
    模型 & BGE 中文嵌入模型、交叉编码器重排模型、可配置大模型 API \\
    客户端 & Chrome 120 及以上，1920$\times$1080 分辨率 \\
    网络 & 校园网或家用宽带，记录上下行带宽和接口时延 \\
    知识库 & 校园公开制度、办事指南、常见问答，按教务、学工、后勤、图书馆分类 \\
    测试集 & 100 道校园问题，四个主题各 25 道，人工标注标准答案和来源 \\
    \bottomrule
  \end{tabular}
\end{table}

测试分为四类。第一，功能测试。按照需求清单逐项验证登录、上传、解析、问答、引用、历史、反馈、统计和配置等功能，检查正常流程和异常流程是否返回预期结果。第二，性能测试。使用脚本模拟并发请求，记录平均响应时间、95 分位响应时间、错误率和资源占用；问答接口的时延拆分为检索、重排、生成三段，分别统计。第三，效果测试。在标注好的测试集上运行不同方案，计算检索命中率、答案正确率、引用正确率和拒答准确率。第四，安全测试。检查未登录访问、越权访问、令牌过期、恶意文件上传、超长输入和频繁请求等场景，确认系统返回正确的错误码且不泄露敏感信息。

功能测试采用“用例—操作—预期结果—实际结果”的记录方式。表~\ref{tab:test-case} 给出主要测试用例的设计，实际执行时逐条填写通过与否、实际输出、耗时和截图编号。对于失败用例，记录复现步骤、日志请求编号和初步原因，修复后回归测试。异常场景与正常场景同等重要，例如检索为空时是否拒答、模型超时是否降级、引用越界是否隐藏，这些行为直接关系到系统是否可信。

\begin{table}[H]
  \centering
  \caption{主要功能测试用例设计}
  \label{tab:test-case}
  \zihao{-5}
  \renewcommand{\arraystretch}{1.05}
  \begin{tabular}{p{1.1cm}p{3.0cm}p{5.5cm}p{3.3cm}}
    \toprule
    编号 & 测试项 & 操作或输入 & 预期结果 \\
    \midrule
    T-01 & 正常登录 & 输入正确账号和密码 & 返回令牌，进入问答页 \\
    T-02 & 密码错误 & 输入错误密码 & 提示账号或密码错误，不签发令牌 \\
    T-03 & 越权访问 & 学生令牌调用资料删除接口 & 返回 403，不执行删除 \\
    T-04 & 上传 PDF & 上传一份可解析的 PDF & 返回文档编号，状态转为成功 \\
    T-05 & 上传非法文件 & 上传可执行文件或超大文件 & 拒绝上传，提示格式或大小错误 \\
    T-06 & 解析失败 & 上传损坏的 DOCX & 状态为失败，展示原因并可重试 \\
    T-07 & 正常问答 & 提问“补办学生证需要什么材料” & 返回答案、引用编号和来源卡片 \\
    T-08 & 证据不足 & 提问知识库中没有的内容 & 拒答，提示换一种问法 \\
    T-09 & 引用校验 & 构造模型返回越界引用 & 隐藏无效引用，降级为摘录回答 \\
    T-10 & 多轮追问 & 先问助学金条件，再问“需要哪些材料” & 改写后检索到助学金材料，答案连贯 \\
    T-11 & 历史查询 & 打开历史页并按关键词筛选 & 只显示当前用户的会话 \\
    T-12 & 统计与配置 & 管理员查看高频问题并修改 $k_1$ & 图表更新，配置保存并校验范围 \\
    \bottomrule
  \end{tabular}
\end{table}

性能测试使用 Locust 或 JMeter 编写脚本，分别设置 1、5、10、20 个并发用户，持续运行 5 分钟，记录请求成功率、平均时延、95 分位时延和吞吐量。由于大模型调用受外部服务影响，性能测试分为“仅检索”和“完整问答”两组：仅检索组不调用大模型，用于评估本地检索链路的性能；完整问答组调用大模型，用于评估端到端体验。两组都记录失败请求的错误码和超时比例。测试期间监控 CPU、内存和网络占用，避免把资源不足导致的超时误判为代码问题。

\section{问答测试集与评价指标}
\label{sec:metrics}

测试集从校园常见问题中整理，覆盖教务管理、校园生活、实验室使用和办事流程四个主题，每个主题 25 道，共 100 道。题目来源包括任务书和开题报告中的场景、学校公开的办事指南、学生访谈中提出的问题，以及作者在日常学习中遇到的实际问题。每道题记录问题、标准答案、来源文档、答案要点和标注依据；对于有多种合理表述的题目，标准答案给出要点集合，而不是唯一字符串。测试集构成如表~\ref{tab:testset} 所示。

\begin{table}[H]
  \centering
  \caption{问答测试集构成}
  \label{tab:testset}
  \zihao{5}
  \renewcommand{\arraystretch}{1.25}
  \begin{tabular}{p{2.6cm}p{1.4cm}p{8.4cm}}
    \toprule
    主题 & 题量 & 示例 \\
    \midrule
    教务管理 & 25 & 补办学生证、成绩单打印、转专业申请、考试违纪处理 \\
    校园生活 & 25 & 宿舍报修、校园网连接、电费充值、心理咨询预约 \\
    实验室使用 & 25 & 实验室预约、设备借用、安全考试、开放时间 \\
    办事流程 & 25 & 在读证明、党员材料、档案转递、离校手续 \\
    \bottomrule
  \end{tabular}
\end{table}

评价指标包括检索命中率、答案正确率、引用正确率、拒答准确率和平均响应时延。检索命中率衡量前 $k$ 条候选是否包含标注来源片段，定义为
\begin{equation}
  \mathrm{HitRate@}k=\frac{1}{N}\sum_{i=1}^{N}\mathbb{I}\left(\mathcal{C}_i(k)\cap\mathcal{G}_i\neq\varnothing\right),
  \label{eq:hitrate}
\end{equation}
其中 $N$ 为题目数量，$\mathcal{C}_i(k)$ 为第 $i$ 题的前 $k$ 条候选片段集合，$\mathcal{G}_i$ 为该题的标准来源片段集合，$\mathbb{I}(\cdot)$ 为指示函数。答案正确率按人工判定，答案覆盖标准答案的全部要点记为正确，覆盖部分要点记为部分正确：
\begin{equation}
  \mathrm{Acc}=\frac{N_{\text{correct}}+0.5N_{\text{partial}}}{N}.
  \label{eq:accuracy}
\end{equation}
引用正确率衡量答案中的引用是否真实支持对应结论，定义为
\begin{equation}
  \mathrm{CiteAcc}=\frac{N_{\text{valid citation}}}{N_{\text{total citation}}},
  \label{eq:citeacc}
\end{equation}
其中有效引用指引用编号存在、片段可访问、且片段内容与所支持结论一致。拒答准确率分别统计“无答案题目正确拒答”和“有答案题目未误拒”的比例，用两者的平均值表示：
\begin{equation}
  \mathrm{RefusalAcc}=\frac{1}{2}\left(\frac{N_{\text{refuse correct}}}{N_{\text{unanswerable}}}+
  \frac{N_{\text{answer correct}}}{N_{\text{answerable}}}\right).
  \label{eq:refusal}
\end{equation}
平均响应时延按检索、重排、生成三段分别统计，并计算端到端平均值和 95 分位值。时延只统计成功返回的请求，超时和错误请求单独统计比例。

人工标注采用双人独立标注、分歧讨论的方式。标注前统一答案要点和引用有效的判定标准；标注时只看问题、答案和引用原文，不看系统名称和实验分组，减少主观偏差。对于争议样本，由第三人复核并记录理由。测试集和标注结果保存在版本库中，实验后不再修改，保证不同方案在同一数据上比较。

\section{对照实验设计}
\label{sec:experiment}

对照实验围绕两个问题展开：检索增强是否比纯大模型回答更准确、更可核查；不同的检索策略对效果有什么影响。实验设置五个组，如表~\ref{tab:groups} 所示。A 组为纯大模型直答，不提供任何校园资料；B 组为 RAG 问答，按检索策略分为 B1（仅向量检索）、B2（BM25 与向量混合）、B3（混合检索加重排序）。A 组用于观察无外部知识时的表现，B 组内部用于比较检索策略的贡献。

\begin{table}[H]
  \centering
  \caption{对照实验分组}
  \label{tab:groups}
  \zihao{5}
  \renewcommand{\arraystretch}{1.25}
  \begin{tabular}{p{1.2cm}p{4.4cm}p{6.8cm}}
    \toprule
    组别 & 检索与生成方式 & 观察目的 \\
    \midrule
    A & 纯大模型直答，不提供证据 & 无校本知识时的准确性与引用能力 \\
    B1 & 仅向量召回，Top-$k$ 后直接生成 & 稠密检索的基线效果 \\
    B2 & BM25 与向量双路召回，RRF 融合 & 混合检索是否提高召回质量 \\
    B3 & 混合检索后交叉编码器重排 & 重排序是否提高顶部证据相关性 \\
    \bottomrule
  \end{tabular}
\end{table}

实验控制以下变量：使用同一套测试题、同一版本的大模型、同一温度和最大生成长度、同一套输出格式要求；B 组使用同一知识库快照，只是检索策略不同；每组运行一次完整测试集，记录原始输出、引用、耗时和错误。为了减少顺序效应，测试题按随机顺序执行；对于可能受缓存影响的指标，实验前清空检索缓存。A 组不提供证据，因此不统计引用正确率，只统计答案正确率和响应时延；B 组统计全部指标。

实验流程分为四步。第一步，构建并冻结知识库，记录文档数量、片段数量和分类分布。第二步，运行 A 组和 B 组，保存每道题的答案、引用、候选片段和耗时。第三步，人工标注答案正确性和引用有效性，标注时不看组别。第四步，汇总指标，按主题和题型分组分析，找出失败案例的共同原因。对于差异明显的指标，进一步做参数敏感性实验，例如改变 $k_1$、$k_2$、$\tau$ 和分块长度，观察指标变化。

\section{结果记录与分析方法}
\label{sec:analysis}

实验结果按“原始记录—汇总指标—失败分析”三层组织。原始记录包括每道题的组别、答案、引用、候选片段、耗时和错误码；汇总指标按组别和主题计算平均值、标准差和样本数；失败分析把错误分为检索未命中、证据不完整、生成偏离证据、引用错误、拒答误判五类，统计各类占比并给出典型案例。这样既能看出整体差异，也能定位问题出在检索还是生成。

分析时需要注意三点。第一，区分“答案正确”和“引用正确”。一个答案可能内容正确但引用错误，也可能引用正确但答案不完整，两者要分开统计。第二，关注拒答的两类错误。把可回答的问题误拒，会降低可用性；把不可回答的问题强行回答，会带来错误信息风险，两类错误的代价不同，需要结合使用场景权衡。第三，避免只报告平均值。不同主题的题目难度差异较大，按主题分组报告可以避免某一类题目拉高或拉低整体指标。

受毕业设计进度限制，完整的对照实验将在系统联调完成后执行。本节给出的测试集、指标、分组和记录方法已经固定，实验执行时只需按表中字段填写实际结果，不需要重新设计评价方案。所有数据以系统实际输出和人工标注为准，不以预期值、经验值或他人系统的结果替代；对于无法复现或超出解释范围的结果，在论文中如实说明。
